\chapter{Digital modulation and demodulation}\label{ch:modulation}

% ------------------------------------------------------------------
\section{Constellations}
A memoryless modulator maps each group of $k$ bits $\vect b=(b_0,\dots,b_{k-1})$ to one of
$M=2^k$ complex points, $s=\mathcal{S}(\vect b)$. The label $\ell=\sum_i b_i2^{k-1-i}$ is
read MSB first and indexes the table of points. LinkStudio builds four families, all normalized to
unit average energy $\E|s|^2=1$:
\begin{align}
  \text{M-PAM:}\quad & s_\ell = \frac{2i-(M-1)}{\sqrt{(M^2-1)/3}},\quad \ell=g(i),\\
  \text{M-QAM:}\quad & s_\ell = \frac{a_{\ell_I}+j\,a_{\ell_Q}}{\sqrt{2(M-1)/3}},\quad
     \ell=(\ell_I,\ell_Q),\ a_\cdot\ \text{Gray }\sqrt M\text{-PAM levels},\\
  \text{M-PSK:}\quad & s_\ell = e^{\,j(2\pi i/M+\phi_0)},\quad \ell=g(i),\quad
     \phi_0=\pi/4\ \text{for QPSK, else }0,
\end{align}
where $g(i)=i\oplus\lfloor i/2\rfloor$ is the binary-reflected Gray code. Adjacent points
therefore differ in exactly one bit. The denominators are the average energies of the
unnormalized PAM and QAM grids, $\frac1M\sum_i(2i-M+1)^2=(M^2-1)/3$ per dimension. Custom
constellations are normalized the same way.

\begin{figure}[htbp]
\centering
\begin{tikzpicture}[scale=0.85]
  \draw[->,black!50] (-4.4,0)--(4.4,0) node[right,lbl]{I};
  \draw[->,black!50] (0,-4.4)--(0,4.4) node[above,lbl]{Q};
  \foreach \x/\bx in {-3/00,-1/01,1/11,3/10}
    \foreach \y/\by in {-3/00,-1/01,1/11,3/10} {
      \fill[accent] (\x,\y) circle (2.5pt);
      \node[font=\scriptsize\ttfamily,above=1pt] at (\x,\y) {\textcolor{accent2}{\bx}\textcolor{accent3}{\by}};
    }
  \foreach \x/\bx in {-3/00,-1/01,1/11,3/10}
    \node[font=\scriptsize\ttfamily,accent2,below] at (\x,-4.0) {\bx};
  \node[lbl,accent2,below] at (0,-4.45) {I bits $(b_0b_1)$};
  \node[lbl,accent3,rotate=90] at (-4.7,0) {Q bits $(b_2b_3)$};
  \draw[dashed,black!25] (-2,-4) -- (-2,4) (0,-4)--(0,4) (2,-4)--(2,4) (-4,-2)--(4,-2) (-4,2)--(4,2);
\end{tikzpicture}
\caption{16-QAM with the labels produced by \code{qam(16)}. The first half of the label selects
the in-phase level and the second half the quadrature level, each with a Gray 4-PAM labelling.
Dashed lines are the decision boundaries.}
\label{fig:qam16}
\end{figure}

\codeexcerpt{qam_build}{Square QAM as the product of two Gray PAMs (\file{dsp/constellations.py}).}

\begin{incode}
\blk{mapper} (\file{phy/modulation.py}) builds the constellation in \code{setup()} and maps
frames of $n k$ bits to $n$ symbols. It sets \meta{constellation}, \meta{bits_per_symbol},
\meta{power}$=1$, \meta{sps}$=1$ and the symbol rate \prm{symbol_rate} as \code{fs}.
\blk{demapper} takes the constellation from the metadata unless one is given explicitly.
\end{incode}

% ------------------------------------------------------------------
\section{Optimal detection}
For the observation $y=s+w$ with $w\sim\CN(0,\sigma^2)$ and equiprobable symbols, the
maximum-likelihood decision is the nearest point:
\begin{equation}
  \hat s=\argmax_{s\in\mathcal{S}} p(y\,|\,s)=\argmax_{s}\frac{1}{\pi\sigma^2}e^{-|y-s|^2/\sigma^2}
  =\argmin_{s\in\mathcal{S}}|y-s|^2 .
\end{equation}
\code{Constellation.nearest} evaluates all $M$ distances in chunks of $2^{16}$ symbols, which
bounds memory. \prm{method}\code{=hard} returns the label bits of the nearest point.

\subsection{Bit log-likelihood ratios}\label{sec:llr}
Soft-input decoders need, for every bit, the a-posteriori log-likelihood ratio. With
$\mathcal{S}_i^{(b)}$ the subset of points whose $i$-th label bit equals $b$, and equiprobable
bits,
\begin{equation}
  L(b_i)=\ln\frac{P(b_i=0\,|\,y)}{P(b_i=1\,|\,y)}
  =\ln\!\sum_{s\in\mathcal{S}_i^{(0)}}\!e^{-|y-s|^2/\sigma^2}
   -\ln\!\sum_{s\in\mathcal{S}_i^{(1)}}\!e^{-|y-s|^2/\sigma^2}.
  \label{eq:llr-exact}
\end{equation}
Keeping only the dominant term of each sum gives the max-log approximation
\begin{equation}
  L(b_i)\approx\frac{1}{\sigma^2}\Big(\min_{s\in\mathcal{S}_i^{(1)}}|y-s|^2-\min_{s\in\mathcal{S}_i^{(0)}}|y-s|^2\Big),
  \label{eq:llr-maxlog}
\end{equation}
which is exact for BPSK and loses only a few tenths of a dB with good codes. A positive LLR
means bit $0$ throughout LinkStudio. The variance $\sigma^2$ may differ from symbol to symbol, for
example after zero-forcing equalization of a fading channel; \meta{noise_var} may therefore be
an array broadcastable to the symbols.
\codeexcerpt{demap_soft}{Exact (log-sum-exp) and max-log LLRs.}

\begin{keypoint}
The LLR scale matters. Equalizers update \meta{noise_var} with their noise enhancement
($\sigma^2/|h|^2$ for a one-tap ZF equalizer, $\sigma^2\|\vect w\|^2$ for an FIR equalizer,
the post-detection variance for MIMO). The decoder downstream then receives LLRs with the
correct reliability without any manual tuning.
\end{keypoint}

% ------------------------------------------------------------------
\section{Error probability over AWGN}\label{sec:ber-awgn}
With $\Q(x)=\frac12\erfc(x/\sqrt2)$ the Gaussian tail, the classic results implemented in
\file{theory.py} are:
\begin{itemize}
  \item \textbf{BPSK and Gray QPSK}: $P_b=\Q\big(\sqrt{2E_b/N_0}\big)$. QPSK is two orthogonal BPSK
  channels sharing the symbol energy.
  \item \textbf{Gray M-PAM}, exact (Cho and Yoon~\cite{cho2002}): with $L=M$, $m=\log_2L$ and
  $\gamma_b=E_b/N_0$,
  \begin{equation}
    P_b=\frac1m\sum_{q=1}^{m}\frac1L\!\!\sum_{i=0}^{(1-2^{-q})L-1}\!\!
    (-1)^{\lfloor i2^{q-1}/L\rfloor}\Big(2^{q-1}-\Big\lfloor\frac{i2^{q-1}}{L}+\frac12\Big\rfloor\Big)
    \erfc\!\Big((2i+1)\sqrt{\tfrac{3m\gamma_b}{L^2-1}}\Big).
  \end{equation}
  \item \textbf{Square Gray M-QAM}: the in-phase and quadrature components are independent
  $\sqrt M$-PAM signals carrying half the bits each, so $P_b^{\text{QAM}}(M)=P_b^{\text{PAM}}(\sqrt M)$
  at the same $E_b/N_0$.
  \item \textbf{M-PSK}: the symbol error rate is Craig's integral~\cite{craig1991}
  \begin{equation}
    P_s=\frac1\pi\int_0^{(M-1)\pi/M}\exp\!\Big(-\frac{E_s}{N_0}\frac{\sin^2(\pi/M)}{\sin^2\theta}\Big)d\theta,
    \qquad P_b\approx P_s/k ,
  \end{equation}
  where the Gray approximation $P_b\approx P_s/k$ is tight at moderate SNR.
\end{itemize}

\begin{figure}[htbp]
\centering
\begin{tikzpicture}
\begin{semilogyaxis}[berplot,width=12cm,height=8cm,xmin=-2,xmax=16,ymin=1e-6,ymax=0.5,
  legend columns=2,legend pos=south west,legend style={font=\scriptsize}]
  \addplot[th=accent] table[x=ebn0,y=bpsk]{\datafile{ber_awgn_theory}};
  \addplot[sim=accent] table[x=ebn0,y=bpsk]{\datafile{ber_awgn_sim}};
  \addplot[th=accent3] table[x=ebn0,y=psk8]{\datafile{ber_awgn_theory}};
  \addplot[sim={accent3,mark=square*}] table[x=ebn0,y=psk8]{\datafile{ber_awgn_sim}};
  \addplot[th=accent2] table[x=ebn0,y=qam16]{\datafile{ber_awgn_theory}};
  \addplot[sim={accent2,mark=triangle*}] table[x=ebn0,y=qam16]{\datafile{ber_awgn_sim}};
  \addplot[th=accent4] table[x=ebn0,y=qam64]{\datafile{ber_awgn_theory}};
  \addplot[sim={accent4,mark=diamond*}] table[x=ebn0,y=qam64]{\datafile{ber_awgn_sim}};
  \draw[black!60,dashed] (axis cs:-1.59,1e-6) -- (axis cs:-1.59,0.5)
     node[pos=0.08,right,font=\scriptsize,black!70]{Shannon limit $-1.59$\,dB};
  \legend{BPSK/QPSK theory,simulated,8-PSK theory,simulated,16-QAM theory,simulated,64-QAM theory,simulated}
\end{semilogyaxis}
\end{tikzpicture}
\caption{Bit error rate over AWGN: simulated (markers, chain
\blk{random_bits}$\to$\blk{mapper}$\to$\blk{awgn}$\to$\blk{demapper}$\to$\blk{ber_meter}, up to
200 errors per point) and exact theory (lines). Points with fewer than 10 errors are omitted.}
\label{fig:ber-awgn}
\end{figure}

\paragraph{The ultimate limit.}
For a spectral efficiency $\eta$ bit/s/Hz, the AWGN capacity $C=\log_2(1+\mathrm{SNR})$ with
$\mathrm{SNR}=\eta E_b/N_0$ requires $E_b/N_0\ge(2^\eta-1)/\eta$, which tends to
$\ln2=-1.59$\,dB as $\eta\to0$ (\code{ebn0_min_db}). The gap between this limit and the uncoded
curves of \cref{fig:ber-awgn} is what channel coding recovers (\cref{ch:coding}).
